
\subsection*{Remarques}

Cette section présente les principales difficultés rencontrées au cours du projet, ainsi que les solutions mises en place pour y répondre.

Une première difficulté rencontrée lors du projet concernait la gestion de la rotation des tuiles dans l'achievement. Dans un premier temps, les rotations étaient définies séparément dans un fichier dédié, avec des fonctions permettant de consulter et de modifier la rotation d’une tuile.\\
Cette approche compliquait toutefois la logique du jeu, car elle nécessitait de maintenir deux tableaux distincts : un pour les tuiles et un autre pour leurs rotations. Ces tableaux devaient ensuite être initialisés et synchronisés dans le code principal, ce qui rendait la boucle de jeu plus lourde et plus difficile à maintenir.\\
Pour simplifier cette gestion, une seconde solution a été mise en place et qui consiste à définir une nouvelle structure de tuile regroupant à la fois la tuile et sa rotation. Cette solution permet de manipuler une tuile comme une seule entité, ce qui facilite son utilisation dans la boucle de jeu et rend le code plus clair.

C'était donc la raison pour laquelle on a un ordre d'achievements un peu bizarre et donc le bon ordre des achievements est :
\begin{enumerate}
    \item Achiev0
    \item Achiev1.2
    \item Achiev1.6
    \item Achiev2
    \item Achiev3
\end{enumerate}

Concernant le second achievement la fonction qui était un peu difficile à implémenter est la fonction récursive cherchant le chemin vers la tuile-objective, l’objectif principale était d’éviter une complexité trop élevée. Après réflexion, cette partie s’est révélée moins difficile que prévu, la solution étant surtout intuitive et naturelle à mettre en place.

Le troisième achievement était simple pour l'implémenter puis inclure les fonctions dédiés dans la boucle du jeu, sauf qu'on avait un problème concernant les ressources mais après utilisation de la fonction \texttt{stash\_to\_string} pour visualiser la progression des ressources on a conclu que la faute était tout simplement à cause d'une boucle.

\subsection*{Conclusion}

Ce projet a permit de mettre en pratique les notions de programmation en langage C vues au cours du semestre, à travers la réalisation d’un jeu complet et fonctionnel. La conception modulaire du code a facilité la compréhension du projet et son évolution progressive.

Les différentes difficultés rencontrées ont été l’occasion de réfléchir à des choix d’implémentation plus adaptés et d’améliorer la structure du programme. Au final, le jeu obtenu respecte les objectifs du projet et constitue une base claire pour d’éventuelles améliorations.